\documentclass[a4paper]{article}

\usepackage{graphicx}

\usepackage{amsmath,amssymb}

\usepackage{minted}

\usepackage{multirow}

\usepackage{float}

\usepackage{algorithm}

\usepackage{algpseudocode}

\usepackage{todonotes}

\usepackage{forest}

\usepackage{xstring}

\usepackage{xcolor}

\usepackage[most]{tcolorbox}

\usepackage{xparse}

\usepackage{natbib}

\usepackage{tikz}

\usetikzlibrary{graphs,graphdrawing}

\usegdlibrary{layered}

% for external graphviz

\input{/home/donkarlo/Dropbox/repo/nd_language_ai_project/src/nd_language_ai/formal/computer/programming/tex/latex/code_clips/external_graphviz.tex}

% for tags

\input{/home/donkarlo/Dropbox/repo/nd_language_ai_project/src/nd_language_ai/formal/computer/programming/tex/latex/parts/control_sequence/kind/semtag/semtag.tex}

% for links

\input{/home/donkarlo/Dropbox/repo/nd_language_ai_project/src/nd_language_ai/formal/computer/programming/tex/latex/code_clips/seehere.tex}

\title{Bayesian Periodicity Detection}

\date{}

\author{Mohammad Rahmani}

\begin{document}

\maketitle

% ND_MIND_MIRROR_CONTENT

\section{Purpose}

Bayesian periodicity detection provides a probabilistic alternative to a purely similarity-based motif score. Given a candidate time-series segment, the goal is to quantify the evidence that the observations were generated by a periodic process and, conditional on periodicity, to infer which periods are plausible. The Bayesian formulation is particularly useful when the desired output is a probability-like confidence rather than only a distance, correlation, or periodogram peak. Gregory and Loredo introduced a Bayesian model-comparison framework for periodic signals with unknown period and unknown waveform \citep{gregory-1992-a-new-method-for-the-detection-of-a-periodic-signal-of-unknown-shape-and-period}. Gregory later extended the framework to nonuniformly sampled measurements with independent Gaussian noise \citep{gregory-1999-bayesian-periodic-signal-detection-i-analysis-of-20-years-of-radio-flux-measurements-of-the-x-ray-binary-ls-i-plus-61-303}.

Let the candidate data be

\begin{equation}
D = \{(t_i,y_i,\sigma_i)\}_{i=1}^{N},
\end{equation}

where $t_i$ is the observation time, $y_i$ is the measured value, and $\sigma_i$ is an optional measurement uncertainty.

\section{Competing hypotheses}

A minimal Bayesian test compares two hypotheses:

\begin{itemize}

    \item $H_C$: the candidate is described by a constant or nonperiodic reference model plus noise.

    \item $H_P$: the candidate contains a periodic modulation with an unknown period and, in the Gregory--Loredo formulation, an unknown waveform.

\end{itemize}

For a stricter test, a third hypothesis can be included:

\begin{itemize}

    \item $H_{NP}$: the signal is genuinely variable, but the variability is nonperiodic.

\end{itemize}

The three-hypothesis comparison is useful because it prevents arbitrary variability from being interpreted automatically as periodicity. This distinction is explicitly used in the Gaussian-data extension of the Gregory--Loredo framework \citep{gregory-1999-bayesian-periodic-signal-detection-i-analysis-of-20-years-of-radio-flux-measurements-of-the-x-ray-binary-ls-i-plus-61-303}.

\section{Bayesian model evidence}

For any hypothesis $H_k$ with parameters $\theta_k$, the Bayesian evidence is

\begin{equation}
Z_k
=
p(D\mid H_k)
=
\int p(D\mid \theta_k,H_k)\,p(\theta_k\mid H_k)\,d\theta_k.
\end{equation}

The integral averages the likelihood over the allowed parameter values instead of selecting only the best-fitting parameters. Therefore, a very flexible periodic model is automatically penalized when most of its parameter space does not explain the data well.

The Bayes factor comparing the periodic and constant hypotheses is

\begin{equation}
B_{PC}
=
\frac{p(D\mid H_P)}{p(D\mid H_C)}
=
\frac{Z_P}{Z_C}.
\end{equation}

A Bayes factor is an evidence ratio, not by itself a posterior probability. If the prior model probabilities are $\pi_P=P(H_P)$ and $\pi_C=P(H_C)$, and these are the only two hypotheses, then

\begin{equation}
P(H_P\mid D)
=
\frac{Z_P\pi_P}{Z_P\pi_P+Z_C\pi_C}
=
\frac{B_{PC}\pi_P}{B_{PC}\pi_P+\pi_C}.
\end{equation}

For equal prior model probabilities, $\pi_P=\pi_C$, this simplifies to

\begin{equation}
P(H_P\mid D)
=
\frac{B_{PC}}{1+B_{PC}}.
\end{equation}

This quantity can be used as the probability that the candidate is periodic, provided that the stated model set and model priors are accepted.

If $H_C$, $H_P$, and $H_{NP}$ are compared simultaneously, then

\begin{equation}
P(H_P\mid D)
=
\frac{Z_P\pi_P}
{Z_C\pi_C+Z_P\pi_P+Z_{NP}\pi_{NP}}.
\end{equation}

\section{Posterior probability of the period}

The probability that a signal is periodic and the probability that it has a particular period are different quantities. Conditional on $H_P$, the posterior density of the period $T$ is

\begin{equation}
p(T\mid D,H_P)
=
\frac{p(D\mid T,H_P)p(T\mid H_P)}
{p(D\mid H_P)},
\end{equation}

where the period-specific marginal likelihood is

\begin{equation}
p(D\mid T,H_P)
=
\int
p(D\mid T,\eta,H_P)
\,p(\eta\mid T,H_P)
\,d\eta,
\end{equation}

and $\eta$ denotes nuisance parameters such as phase, waveform parameters, amplitude, and noise parameters.

Suppose a motif-discovery stage proposes a candidate period $T^{\ast}$. The posterior probability that the true period lies in a tolerance interval around this candidate is

\begin{equation}
P\left(\left|T-T^{\ast}\right|<\Delta T\mid D,H_P\right)
=
\int_{T^{\ast}-\Delta T}^{T^{\ast}+\Delta T}
p(T\mid D,H_P)\,dT.
\end{equation}

Therefore, two useful outputs are available:

\begin{equation}
S_{\mathrm{periodic}}
=
P(H_P\mid D),
\end{equation}

which quantifies how probable the periodic hypothesis is, and

\begin{equation}
S_{T^{\ast}}
=
P\left(\left|T-T^{\ast}\right|<\Delta T\mid D,H_P\right),
\end{equation}

which quantifies how plausible a particular candidate period is, conditional on the signal being periodic.

\section{Gaussian time-series likelihood}

For continuously valued sensory time series, a simple observation model is

\begin{equation}
y_i
=
g\!\left(\phi_i;\beta\right)+\epsilon_i,
\end{equation}

with

\begin{equation}
\phi_i
=
\operatorname{frac}\!\left(\frac{t_i-t_0}{T}\right),
\end{equation}

and

\begin{equation}
\epsilon_i\sim\mathcal{N}(0,s_i^2).
\end{equation}

Here $g(\phi_i;\beta)$ is the periodic waveform evaluated at phase $\phi_i$. The corresponding likelihood is

\begin{equation}
p(D\mid T,\beta,H_P)
=
\prod_{i=1}^{N}
\frac{1}{\sqrt{2\pi s_i^2}}
\exp\left[
-\frac{\left(y_i-g(\phi_i;\beta)\right)^2}{2s_i^2}
\right].
\end{equation}

The noise variance may be known, estimated, or marginalized as a nuisance parameter. The Gaussian extension of the Gregory--Loredo method was designed for this type of sampled time-series setting and allows nonuniform observation times \citep{gregory-1999-bayesian-periodic-signal-detection-i-analysis-of-20-years-of-radio-flux-measurements-of-the-x-ray-binary-ls-i-plus-61-303}.

\section{Gregory--Loredo unknown-waveform model}

The Gregory--Loredo method does not require the periodic waveform to be sinusoidal \citep{gregory-1992-a-new-method-for-the-detection-of-a-periodic-signal-of-unknown-shape-and-period}. For a trial period $T$, the observations are folded into phase and the phase interval is divided into $m$ bins. Each bin has its own signal level, so a piecewise-constant periodic waveform can approximate many different cycle shapes.

A periodic model can therefore be written conceptually as

\begin{equation}
g(\phi)
=
a_j,
\qquad
\phi\in I_j,
\qquad
j=1,\ldots,m,
\end{equation}

where $I_j$ is the $j$th phase bin and $a_j$ is its signal level. Instead of fixing a single number of bins, phase offset, waveform, or noise value, the Bayesian calculation marginalizes over them. The resulting evidence rewards periods that consistently align repeated observations while penalizing unnecessarily complicated periodic explanations.

For a sensory sequence, the practical search variables are therefore

\begin{equation}
T,\quad m,\quad \varphi_0,\quad \{a_j\}_{j=1}^{m},\quad \text{and noise parameters}.
\end{equation}

The output is not only a best period. It is a posterior distribution over the period together with an evidence measure for periodicity.

\section{Procedure for a candidate subsequence}

\begin{algorithm}[H]
\caption{Bayesian periodicity detection for a candidate subsequence}
\begin{algorithmic}[1]

\Require Candidate data $D$, period range $[T_{\min},T_{\max}]$, model priors $\pi_C,\pi_P$ and optionally $\pi_{NP}$, tolerance $\Delta T$

\Ensure Posterior periodicity probability, period posterior, candidate-period probability, and Bayes factor

\State Specify the constant/reference model $H_C$ and the periodic model $H_P$.

\State Compute the evidence $Z_C=p(D\mid H_C)$ by marginalizing the parameters of $H_C$.

\For{each candidate period $T\in[T_{\min},T_{\max}]$}

    \State Fold the observations into phase using $T$.

    \State Evaluate periodic models over the allowed waveform complexity, phase offset, and noise parameters.

    \State Marginalize nuisance parameters to obtain $p(D\mid T,H_P)$.

\EndFor

\State Normalize over $T$ to obtain $p(T\mid D,H_P)$ and integrate over $T$ to obtain $Z_P$.

\If{$H_{NP}$ is included}

    \State Compute $Z_{NP}=p(D\mid H_{NP})$.

\EndIf

\State Compute $P(H_P\mid D)$ from the evidences and prior model probabilities.

\State Compute $B_{PC}=Z_P/Z_C$.

\State Find $T_{\mathrm{MAP}}=\arg\max_T p(T\mid D,H_P)$ and a posterior credible interval.

\State For a proposed period $T^{\ast}$, integrate $p(T\mid D,H_P)$ over $[T^{\ast}-\Delta T,T^{\ast}+\Delta T]$.

\State \Return $P(H_P\mid D)$, $B_{PC}$, $p(T\mid D,H_P)$, $T_{\mathrm{MAP}}$, credible interval, and $P(|T-T^{\ast}|<\Delta T\mid D,H_P)$.

\end{algorithmic}
\end{algorithm}

\section{Relation to motif discovery}

Motif discovery and Bayesian periodicity detection answer different questions. Motif discovery asks whether similar subsequences recur, whereas Bayesian periodicity detection asks whether the timing of the observations is better explained by a periodic generative model than by competing models. A repeated motif can occur irregularly and therefore need not be periodic. Conversely, a periodic process may have a nonsinusoidal waveform, which is why an unknown-waveform periodic model is useful.

A suitable combined pipeline is

\begin{equation}
\text{motif discovery}
\longrightarrow
\text{candidate cycle or candidate period}
\longrightarrow
\text{Bayesian periodicity detection}
\longrightarrow
\left\{P(H_P\mid D),\;p(T\mid D,H_P)\right\}.
\end{equation}

In this arrangement, motif discovery acts as a proposal mechanism and the Bayesian model supplies a probabilistic validation stage.

\section{Interpretation for learned sensory cycles}

Assume that training data contain repeated sensory cycles and that a candidate subsequence is later observed. A high value of $P(H_P\mid D)$ indicates that the candidate is better supported as periodic than by the chosen alternatives. A concentrated posterior $p(T\mid D,H_P)$ indicates that the period itself is well identified. If the learned cycle is expected near $T^{\ast}$, then $P(|T-T^{\ast}|<\Delta T\mid D,H_P)$ can be used as the probability mass assigned to that expected period interval.

These probabilities are conditional on the selected hypotheses, priors, likelihood, and period range. They should therefore be interpreted as Bayesian posterior probabilities under the specified model, not as model-free probabilities that a sequence is ``truly periodic.''

\section{Reference summaries}

\subsection{Gregory and Loredo (1992)}

\begin{itemize}

    \item Gregory and Loredo introduce a Bayesian method for detecting a periodic signal when both its period and waveform are unknown \citep{gregory-1992-a-new-method-for-the-detection-of-a-periodic-signal-of-unknown-shape-and-period}.

    \item The method compares periodic and nonperiodic explanations through Bayesian model probabilities rather than relying only on a periodogram peak \citep{gregory-1992-a-new-method-for-the-detection-of-a-periodic-signal-of-unknown-shape-and-period}.

    \item The periodic waveform is represented flexibly by phase bins, allowing nonsinusoidal cycle shapes to be considered \citep{gregory-1992-a-new-method-for-the-detection-of-a-periodic-signal-of-unknown-shape-and-period}.

    \item Period, phase, and waveform-related nuisance parameters are handled probabilistically by marginalization \citep{gregory-1992-a-new-method-for-the-detection-of-a-periodic-signal-of-unknown-shape-and-period}.

    \item The framework provides a principled basis for assigning Bayesian evidence to candidate periodic structure in observed data \citep{gregory-1992-a-new-method-for-the-detection-of-a-periodic-signal-of-unknown-shape-and-period}.

\end{itemize}

\subsection{Gregory (1999)}

\begin{itemize}

    \item Gregory extends the Bayesian periodic-signal framework to nonuniformly sampled data with independent Gaussian noise \citep{gregory-1999-bayesian-periodic-signal-detection-i-analysis-of-20-years-of-radio-flux-measurements-of-the-x-ray-binary-ls-i-plus-61-303}.

    \item The analysis compares constant, periodic, and nonperiodic hypotheses so that variability alone is not sufficient evidence for periodicity \citep{gregory-1999-bayesian-periodic-signal-detection-i-analysis-of-20-years-of-radio-flux-measurements-of-the-x-ray-binary-ls-i-plus-61-303}.

    \item The method derives a posterior probability density for the modulation period instead of reporting only a single best-fitting period \citep{gregory-1999-bayesian-periodic-signal-detection-i-analysis-of-20-years-of-radio-flux-measurements-of-the-x-ray-binary-ls-i-plus-61-303}.

    \item Unknown noise and waveform parameters are treated as nuisance quantities and integrated within the Bayesian calculation \citep{gregory-1999-bayesian-periodic-signal-detection-i-analysis-of-20-years-of-radio-flux-measurements-of-the-x-ray-binary-ls-i-plus-61-303}.

    \item This extension makes the approach directly relevant to continuously valued and irregularly sampled sensory time series \citep{gregory-1999-bayesian-periodic-signal-detection-i-analysis-of-20-years-of-radio-flux-measurements-of-the-x-ray-binary-ls-i-plus-61-303}.

\end{itemize}

\bibliographystyle{plainnat}

\bibliography{/home/donkarlo/Dropbox/repo/nd_shared_research/bibliography/bibliography.bib}

\end{document}
